\documentclass{article}
\usepackage[UTF8]{ctex}
\usepackage{amsmath,amssymb,amsfonts}
\usepackage{enumitem}
\usepackage{geometry}
\geometry{a4paper, margin=1in}
% === 错题本排版优化（不参与题目标识） ===
\setlength{\parindent}{0pt}
\setlength{\parskip}{0.35em plus 0.10em minus 0.08em}
\setlength{\emergencystretch}{3em}
\setlist{itemsep=0.25em, topsep=0.40em, parsep=0pt}
\allowdisplaybreaks[2]
% === 错题本排版优化结束 ===
\begin{document}
\section*{南昌大学2024年数学分析真题}
\section*{119．南昌大学2024年数学分析试卷}

\subsection{题目1}
一．（10 分）求极限 \(\lim _{n \rightarrow \infty} \frac{\sqrt[n]{2024}\left(1-\cos \frac{1}{n^{2}}\right) n^{3}}{\sqrt{n^{2}+1}-n}\) ．

\subsection{题目2}
二．（10 分）求定积分 \(\int_{0}^{\pi} e^{x} \cos ^{2} x \mathrm{~d} x\) ．

\subsection{题目3}
三．（10 分）计算曲线积分

\[
I=\int_{C} \frac{y^{2}}{\sqrt{1+x^{2}}} \mathrm{~d} x+\left[4 x+2 y \ln \left(x+\sqrt{1+x^{2}}\right)\right] \mathrm{d} y
\]

其中 \(C\) 是圆周 \(x^{2}+y^{2}=1\) 由点 \(A(1,0)\) 依逆时针方向到点 \(B(-1,0)\) 的半圆．

\subsection{题目4}
四．（10 分）求幂级数 \(\sum_{n=0}^{\infty} \frac{x^{4 n+1}}{4 n+1}\) 的和函数．

\subsection{题目5}
五．（15 分）设函数 \(f(x)\) 在闭区间 \([a, b]\) 上连续，在开区间 \((a, b)\) 内二阶可导，证明：存在 \(\xi \in(a, b)\) ，使得

\[
f(b)-2 f\left(\frac{a+b}{2}\right)+f(a)=\frac{(b-a)^{2}}{4} f^{\prime \prime}(\xi)
\]

\subsection{题目6}
六．（15分）求二元函数

\[
f(x, y)= \begin{cases}\frac{x^{3}-y^{3}}{x^{2}+y^{2}}, & x^{2}+y^{2} \neq 0 \\ 0, & x^{2}+y^{2}=0\end{cases}
\]

的偏导数，并讨论 \(f(x, y)\) 在原点 \((0,0)\) 处的可微性．

\subsection{题目7}
七．（15 分）讨论级数 \(\sum_{n=1}^{\infty} \frac{(-1)^{n}}{n^{\alpha+\frac{1}{n}}}(\alpha>0)\) 的绝对收敛性与条件收敛性．

\subsection{题目8}
八．（15 分）求积分 \(I=\int_{0}^{+\infty} \frac{e^{-x^{2}}-e^{-2 x^{2}}}{x} \mathrm{~d} x\) ．

\subsection{题目9}
九．（15 分）设 \(x_{0}=a \in\left(0, \frac{\pi}{2}\right), x_{n}=\sin x_{n-1}, n=1,2, \cdots\) ．证明：

\begin{itemize}
  \item[（1）] 数列 \(\left\{x_{n}\right\}\) 极限存在且为零。
  \item[（2）] \(\lim _{n \rightarrow \infty} \sqrt{\frac{n}{3}} x_{n}=1\)
\end{itemize}

\subsection{题目10}
十、（15 分）设函数 \(f(x)\) 在区间 \((0,+\infty)\) 内可导，且 \(\lim _{x \rightarrow+\infty} f(x)\) 与 \(\lim _{x \rightarrow+\infty} f^{\prime}(x)\) 均存在，证明：

\[
\lim _{x \rightarrow+\infty} f^{\prime}(x)=0
\]

\subsection{题目11}
十一、（20 分）设 \(f(x)\) 在区间 \([a,+\infty)\) 上一致连续，\(g(x)\) 在区间 \([a,+\infty)\) 上连续，且\(\lim _{x \rightarrow+\infty}[f(x)-g(x)]=0\) ，证明：\(g(x)\) 在区间 \([a,+\infty)\) 上一致连续．
\end{document}
